\documentclass[11pt]{article}
\usepackage[margin=1in]{geometry}
\usepackage{amsmath,amssymb}
\usepackage{booktabs}
\usepackage{array}
\usepackage[hidelinks]{hyperref}
\setlength{\parskip}{0.5em}
\setlength{\parindent}{0pt}

\title{GNN-MPC Cost Function for Square-Rod Forging}
\author{Agility Forge --- control progress notes}
\date{2026-09-26}

\begin{document}
\maketitle

\section{Receding-horizon setup}

At each real hit $t$, the MPC measures the bar's surface displacement
$\mathbf{x}_t \in \mathbb{R}^{N \times 3}$: one displacement vector
$\mathbf{x}_{t,i} = (x_{t,i},\, y_{t,i},\, z_{t,i})$ for each of the
$N \approx 4400$ surface nodes. The component $x$ runs along the bar (clamp to
free end); $y$ and $z$ lie across it, in the plane of a cross-section.

It then plans the next $K = 10$ hits. Hit $k$ has three controls,
\[
  \mathbf{u}_k = (d_k,\; R_k,\; s_k), \qquad k = 1, \dots, K,
\]
where $d_k$ is the station (where the 19.3\,mm die band starts along the bar),
$R_k$ is the angle the bar is turned before the hit, and $s_k$ is the stroke
(how far each die moves in, in mm).

The GNN (message-passing network $f_{\mathrm{GNN}}$) predicts the shape after
each planned hit, starting from the measured state:
\[
  \hat{\mathbf{x}}_0 = \mathbf{x}_t, \qquad
  \hat{\mathbf{x}}_k = \hat{\mathbf{x}}_{k-1} + f_{\mathrm{GNN}}\!\left(\hat{\mathbf{x}}_{k-1},\, \mathbf{u}_k\right),
  \qquad k = 1, \dots, K.
\]
Only the first planned hit $\mathbf{u}_1$ is applied to the simulator. The
simulator returns the real new shape $\mathbf{x}_{t+1}$, and the MPC plans
again.

\section{Optimization problem}

\begin{equation}
  \min_{\mathbf{u}_1, \dots, \mathbf{u}_K} \; J
  \quad \text{subject to} \quad
  \begin{cases}
    d_{\min} \le d_k \le d_{\max} \\
    0^\circ \le R_k \le 360^\circ \\
    0 \le s_k \le s_{\max}
  \end{cases}
  \qquad k = 1, \dots, K,
\end{equation}
with $d_{\min} = 17.7$\,mm (the 800\,$^\circ$C point of the initial temperature
profile), $d_{\max} = 75.3$\,mm, and $s_{\max} = 2$\,mm.

\section{Cost function}

\begin{equation}
  \boxed{\; J \;=\; E \;+\; w\, E_{\perp} \;+\; \frac{E_0}{K} \Big[\, e\, P_{s} \;+\; P_{\Delta} \,\Big] \;}
\end{equation}
where
\begin{itemize}
  \item[(1)] $E$ is the \textbf{error term}: distance of the predicted shape from the target;
  \item[(2)] $E_{\perp}$ is the \textbf{cross-section term}, weighted by $w$;
  \item[(3)] $P_{s}$ is the \textbf{stroke-size penalty}, weighted by $e$;
  \item[(4)] $P_{\Delta}$ is the \textbf{control-change penalty}, with its own weights $\lambda_d, \lambda_R, \lambda_s$.
\end{itemize}
Each term is defined below.

\paragraph{(1) Error term (always used).}
Each surface node's squared distance from its position in the target shape
$\mathbf{x}^{*}$, summed over all nodes and all planned hits:
\begin{equation}
  E \;=\; \sum_{k=1}^{K} \sum_{i=1}^{N} \left\lVert \hat{\mathbf{x}}_{k,i} - \mathbf{x}^{*}_{i} \right\rVert^{2}
  \;=\; \sum_{k=1}^{K} \sum_{i=1}^{N} \Big[ (\hat{x}_{k,i} - x^{*}_{i})^{2} + (\hat{y}_{k,i} - y^{*}_{i})^{2} + (\hat{z}_{k,i} - z^{*}_{i})^{2} \Big].
\end{equation}
For the square target, about 99\% of $E$ at the start comes from the first
(lengthwise) part, because the target bar is 35\,mm longer than the billet.

\paragraph{(2) Cross-section term.}
The same distance, counting only the across-the-bar directions $y$ and $z$:
\begin{equation}
  E_{\perp} \;=\; \sum_{k=1}^{K} \sum_{i=1}^{N} \Big[ (\hat{y}_{k,i} - y^{*}_{i})^{2} + (\hat{z}_{k,i} - z^{*}_{i})^{2} \Big].
\end{equation}
With weight $w$, a sideways miss costs $(1 + w)$ times as much as the same miss
along the bar. This makes the cross-section shape count in the cost.

\paragraph{(3) Stroke-size penalty.}
With the stroke scaled to $[0, 1]$, $\tilde{s}_k = s_k / s_{\max}$:
\begin{equation}
  P_{s} \;=\; \sum_{k=1}^{K} \tilde{s}_k^{\,2}.
\end{equation}

\paragraph{(4) Control-change penalty.}
With the station scaled to $[0, 1]$,
$\tilde{d}_k = (d_k - d_{\min}) / (d_{\max} - d_{\min})$, and $\mathbf{u}_0$ the
last hit actually applied:
\begin{equation}
  P_{\Delta} \;=\; \sum_{k=1}^{K} \Big[
    \lambda_d \,(\tilde{d}_k - \tilde{d}_{k-1})^{2}
    \;+\; \lambda_R \,\sin^{2}\!\left(R_k - R_{k-1}\right)
    \;+\; \lambda_s \,(\tilde{s}_k - \tilde{s}_{k-1})^{2}
  \Big].
\end{equation}
The angle uses $\sin^2$ because the two dies press opposite sides of the bar,
so turning by $0^\circ$ or $180^\circ$ gives the same hit; the penalty is
largest for a $90^\circ$ change.

\paragraph{Scaling of (3) and (4).}
$E_0$ is the value of the error term $E$ at the optimizer's starting guess, so
$E_0 / K$ is one hit's share of the plan's starting error. A penalty weight of
1 therefore means ``a unit penalty on one hit costs as much as one hit's worth
of error'', independent of the error's size ($\sim\!10^{6}$\,mm$^2$ at the
start). The cross-section weight $w$ is not scaled this way; it is in the
error's own units.

\section{Weights used in each experiment}

$M$ is the number of message-passing steps in the GNN that did the planning.

\begin{center}
\small
\renewcommand{\arraystretch}{1.25}
\begin{tabular}{@{}l c c c c c@{}}
\toprule
Experiment & $M$ & $w$ & $e$ & $\lambda_d$ & $\lambda_R,\ \lambda_s$ \\
\midrule
E0: first MPC run (original cost) & 5 & 0 & 0 & 0 & 0 \\
E1: change penalty, all controls & 3 & 0 & 0 & 0.01, 0.1, 1 & 0.01, 0.1, 1 \\
E1: change penalty, station only & 3 & 0 & 0 & 0.01, 0.1, 1 & 0 \\
E2: cross-section / stroke size & 3 & 30, 100, 300 & 0, 0.1, 1 & 0 & 0 \\
E3: refinement grid & 3 & 30, 50, 100 & 0.03, 0.1, 0.3 & 0 & 0 \\
E4: robustness (5 seeds each) & 3 & 30 or 100 & 0.1 or 0.3 & 0 & 0 \\
Real-simulator runs & 3 & 30 or 100 & 0, 0.1 or 0.3 & 0 & 0 \\
\midrule
\textbf{Lead design} & \textbf{3} & \textbf{30} & \textbf{0.3} & \textbf{0} & \textbf{0} \\
\bottomrule
\end{tabular}
\end{center}

With the lead design the cost reduces to
\begin{equation}
  J_{\text{lead}} \;=\; E \;+\; 30\, E_{\perp} \;+\; 0.3\, \frac{E_0}{K} \sum_{k=1}^{K} \left(\frac{s_k}{2\,\text{mm}}\right)^{2}.
\end{equation}

\section{Solver}

$J$ is minimized with SLSQP (sequential quadratic programming with a
quasi-Newton Hessian), using exact gradients $\partial J / \partial \mathbf{u}$
obtained by back-propagating through the $K$ chained GNN predictions. For
numerical conditioning the optimizer works on controls rescaled to $[0,1]$ and
on $J$ divided by its value at the starting guess; this changes the units only,
not the minimizer. The starting guess is the previous plan shifted by one hit
(its last hit repeated).

\end{document}
